\chapter{引言}

\section{研究背景与意义}
电力变压器是电力系统的重要设备。随着设备容量和运行条件的变化，绝缘、冷却和保护装置的状态监测成为运行维护工作的重要组成部分。电机与电气设备的分析方法可参考相关文献\cite{gao1987}。

故障诊断需要将试验数据与设备结构、运行记录及维护经验结合起来。油中溶解气体分析为观察设备内部状态提供信息，但不同测量条件下的结果不能简单地直接比较。

\section{研究内容与论文结构}
本文介绍变压器的基本结构与故障分类，使用相对产气速率说明指标的计算方法，并通过图表展示诊断信息的整理方式。第二章给出方法与示例，第三章总结本文内容并讨论后续工作。

论文中的章、节编号和目录由模板自动生成。图表和公式通过交叉引用关联正文，修改顺序后无需手工重编号。
